\chapter{Conclusão}
\label{chap:conclusoes-e-trabalhos-futuros}

Este trabalho investigou uma forma de aproximar a ideia de autoconstrução dos sistemas vivos ao Aprendizado por Reforço Profundo. Para isso, foi proposta uma configuração experimental que relaciona precariedade perceptiva, curiosidade artificial e manutenção de um corpo artificial. A vida do agente funcionou apenas como uma variável de controle responsável por encerrar o episódio quando seu valor chegava a zero. Ela não foi fornecida à política, não participou do cálculo da recompensa e não informou diretamente o comportamento que deveria ser aprendido. Na condição intrínseca, a coleta de \textit{medikits} também não recebeu recompensa específica. Seu efeito relevante para o aprendizado ocorreu por meio da restauração da percepção e da consequente modificação das observações processadas pelo RND.

Os resultados permitem responder positivamente à questão central dentro do escopo das configurações avaliadas. Nas forças de glaucoma $50$, $100$ e $200$, os agentes treinados apenas com recompensa intrínseca aprenderam políticas capazes de alcançar episódios longos e próximos da duração obtida com recompensa extrínseca. Esse comportamento não ocorreu de maneira consistente na condição sem glaucoma. Com força $0$, a recompensa produzida pelo RND continuou incentivando a exploração, mas permaneceu desacoplada da manutenção da percepção. O desempenho final foi menor, apresentou maior variabilidade entre as sementes e não resultou em uma estratégia clara de coleta na política analisada.

A comparação entre as condições de controle foi essencial para essa interpretação. A ausência de recompensa confirmou que o algoritmo não aprende sistematicamente a prolongar a sobrevivência sem um sinal de orientação. A recompensa extrínseca demonstrou que a arquitetura, o espaço de ações e a tarefa permitem aprender a coleta de recursos quando esse comportamento é valorizado diretamente. A condição intrínseca com glaucoma mostrou que parte desse comportamento também pode emergir quando a recompensa está relacionada às consequências perceptivas da precariedade, em vez de estar diretamente associada aos \textit{medikits}.

Se o único objetivo fosse obter o maior desempenho possível na tarefa, a recompensa extrínseca constituiria uma solução mais simples e direta. A contribuição da proposta está em outro ponto. Nas comparações entre forças de glaucoma, o PPO, o RND, a arquitetura e a definição da recompensa intrínseca permaneceram inalterados. O que mudou foi o corpo perceptivo e a velocidade de sua deterioração. Os comportamentos distintos observados entre as forças indicam que a curiosidade pode ser modulada por meio das condições corporais e ambientais, sem exigir uma nova regra de aprendizado ou uma recompensa específica para cada comportamento desejado.

A relação temporal entre recompensa intrínseca e precariedade forneceu o indício mais direto do funcionamento do mecanismo. Nas forças positivas, o avanço do glaucoma foi acompanhado pela redução da recompensa do RND. A coleta de um recurso restaurou a visão e produziu um aumento abrupto do erro de predição, seguido por nova redução conforme a degradação perceptiva retornava. Na força $0$, a coleta não modificou a percepção e também não produziu uma alteração sistemática na recompensa intrínseca. Dessa forma, os resultados mostram que a relevância dos recursos dependeu da transformação que eles produziram na condição do agente.

As trajetórias e os mapas de Grad-CAM complementaram essa conclusão. As políticas com recompensa apresentaram deslocamentos mais extensos e sensíveis ao ambiente do que as políticas sem recompensa. À medida que a força do glaucoma aumentou, as trajetórias se tornaram menos aderentes à geometria exata dos recursos e mais associadas a estratégias gerais de navegação. Os mapas de Grad-CAM indicaram uma mudança semelhante. Com a visão preservada, os \textit{medikits} puderam receber relevância mais localizada. Sob degradação intensa, paredes, chão, céu, passagens e limites da obstrução assumiram maior importância, pois forneciam referências para o deslocamento antes que a percepção fosse amplamente comprometida.

Em termos conceituais, a precariedade tornou o ambiente significativo para o agente. O \textit{medikit} deixou de ser apenas um elemento visual disponível no cenário e adquiriu significado funcional por restaurar uma condição necessária à continuidade da percepção e da exploração. Esse significado não foi atribuído ao objeto de forma isolada. Ele surgiu da relação entre o recurso, a organização artificial do agente e as consequências da coleta sobre sua experiência futura. O resultado se aproxima da perspectiva enativa adotada como inspiração teórica, segundo a qual a relevância emerge do acoplamento entre agente e ambiente.

As diferenças produzidas pelas forças de glaucoma também são consistentes com os trabalhos de \citeonline{yuri-sem-objetivo} e \citeonline{yuri-ambiente-rico}. Esses estudos mostram que mudanças na energia, na percepção, no corpo e no ambiente alteram as possibilidades de comportamento disponíveis ao agente. Nesta dissertação, a modificação da velocidade de degradação perceptiva produziu diferenças na aprendizagem e na navegação mesmo com a manutenção do mecanismo de curiosidade.

Essa aproximação deve ser compreendida de forma restrita. O agente desenvolvido não é autopoiético, não produz sua própria organização e não determina autonomamente suas condições de viabilidade. O glaucoma, o RND, a arquitetura e a dinâmica de restauração foram definidos pelo projetista. A proposta também não elimina a recompensa do processo de aprendizagem. Sua contribuição consiste em deslocar a orientação do comportamento de uma recompensa que descreve diretamente a solução para um sinal que depende da relação entre condição interna, percepção, história de experiências e ambiente.

\section{Contribuições do Trabalho}
\label{sec:contribuicoes-do-trabalho}

A principal contribuição deste trabalho é a construção de um mecanismo computacional que relaciona precariedade perceptiva e curiosidade artificial. A proposta permite investigar como uma condição corporal pode modificar a percepção e, por consequência, alterar a recompensa intrínseca utilizada pelo aprendizado. Essa formulação oferece uma alternativa experimental à atribuição direta de valor para a ação de coletar recursos.

Essa contribuição representa um deslocamento da engenharia da recompensa para a construção das condições corporais e ambientais. Depois de definido o mecanismo geral de curiosidade, diferentes comportamentos podem ser investigados por meio de alterações no corpo perceptivo sem modificar o PPO ou introduzir recompensas específicas para cada força de glaucoma. O trabalho mostra, assim, uma forma de orientar a curiosidade por meio do corpo do personagem.

Outra contribuição é a organização de uma configuração controlada com condições sem recompensa, com recompensa extrínseca e com recompensa intrínseca. A inclusão da força $0$ como ablação permitiu distinguir os efeitos gerais da curiosidade dos efeitos produzidos por seu acoplamento à precariedade. A avaliação de diferentes forças de glaucoma também mostrou que o mecanismo permanece funcional sob velocidades distintas de degradação, embora intensidades maiores tornem o aprendizado mais lento e a navegação menos precisa.

O trabalho também contribui ao combinar medidas quantitativas e qualitativas. A duração média dos episódios permitiu comparar sobrevivência e estabilidade entre sementes. As trajetórias mostraram como as políticas respondiam a diferentes distribuições espaciais dos recursos. A análise temporal da recompensa permitiu observar o efeito imediato da degradação e da restauração da visão. O Grad-CAM forneceu indícios sobre as regiões utilizadas pelas políticas durante a tomada de decisão. Em conjunto, essas análises reduziram a dependência de uma única métrica e ofereceram uma interpretação mais ampla do comportamento aprendido.

Por fim, os resultados oferecem evidências experimentais de que a curiosidade pode favorecer comportamentos de manutenção quando está vinculada a uma condição que precisa ser preservada. Nas configurações avaliadas, a precariedade transformou a coleta em uma ação relevante sem que uma recompensa específica fosse atribuída ao recurso, à sobrevivência ou à restauração da visão.

Essa contribuição deve ser compreendida como um passo limitado na direção da autonomia artificial. O trabalho não resolve o problema da autonomia nem implementa autopoiese, autonomia constitutiva ou adaptatividade em sentido forte. Ele investiga como uma exigência simplificada de reconstrução funcional pode ser incorporada ao Aprendizado por Reforço e modificar o comportamento aprendido.

\section{Limitações}
\label{sec:limitacoes}

Os resultados estão limitados à tarefa \textit{Health-Gathering} e à dinâmica específica implementada no VizDoom. O ambiente possui uma única sala, espaço de ações reduzido, regras fixas de perda de vida e apenas um tipo principal de recurso. Não é possível afirmar que o mecanismo produziria resultados equivalentes em ambientes mais complexos, com múltiplas necessidades, recursos escassos ou de diferentes naturezas, como alimento, energia, abrigo, ferramentas e itens prejudiciais, além de ameaças móveis ou formas distintas de interação.

O glaucoma artificial representa somente uma forma particular e simplificada de precariedade. A degradação segue uma regra fixa, afeta apenas a percepção e é removida integralmente após a coleta. Outros componentes corporais não se deterioram, não interagem entre si e não exigem regulação simultânea.

A aproximação com a autoconstrução também possui limites claros. O agente não produz seus componentes, não reorganiza sua estrutura, não define sua identidade e não estabelece seus próprios limites de viabilidade. A restauração da visão corresponde a uma reconstrução funcional previamente programada pelo ambiente. Por isso, os resultados não demonstram autopoiese, autonomia constitutiva, adaptatividade em sentido forte ou produção de sentido equivalente à observada em sistemas vivos.

O trabalho também não elimina a recompensa numérica. O RND foi escolhido e implementado pelo projetista, e sua saída continua sendo utilizada pelo PPO para atualizar a política. A contribuição consiste em não atribuir valor direto aos \textit{medikits}, à vida ou à restauração da visão. Ainda assim, o mecanismo pode ser interpretado como uma forma indireta de modelagem da recompensa, pois a estrutura corporal foi projetada de modo a alterar as observações das quais o sinal intrínseco depende.

Outra limitação decorre das escolhas algorítmicas. Os experimentos utilizaram PPO, RND, uma arquitetura convolucional e um conjunto específico de hiperparâmetros. Os resultados não permitem concluir que o mesmo padrão ocorrerá com outros algoritmos, métodos de curiosidade ou arquiteturas. Em particular, modelos recorrentes podem utilizar memória para compensar a perda visual, enquanto outros sinais intrínsecos podem responder de maneira diferente às observações degradadas.

\section{Trabalhos Futuros}
\label{sec:trabalhos-futuros}

Como continuidade, o mecanismo pode ser avaliado em outros ambientes e tarefas que exijam navegação, manipulação de objetos e manutenção de diferentes condições internas. A introdução de necessidades simultâneas, como energia, integridade corporal e temperatura, permitiria investigar como o agente organiza comportamentos quando diferentes dimensões de viabilidade competem entre si.

Também seria relevante comparar o RND com outros métodos de motivação intrínseca já discutidos neste trabalho, como a exploração baseada em pseudocontagens \cite{bellemare2016count} e o \textit{Intrinsic Curiosity Module} (ICM) \cite{icm}. Também poderiam ser avaliadas diferentes representações para o cálculo da curiosidade, como as características aleatórias e aprendidas investigadas por \citeonline{burda2019large}. Essas comparações permitiriam verificar se a relação observada depende especificamente do erro de predição do RND ou se pode ser reproduzida por diferentes medidas de novidade, surpresa e controlabilidade. O uso de arquiteturas recorrentes também poderia esclarecer se a memória permite compensar parcialmente a degradação visual e produzir estratégias mais eficientes sob glaucoma intenso.

Também seria importante comparar diretamente a eficiência da recompensa extrínseca e do mecanismo corporal sob diferentes alterações da tarefa. Essa avaliação poderia verificar se a vantagem da proposta não está apenas em resolver a configuração atual, mas em reutilizar o mesmo mecanismo de curiosidade quando o corpo ou o ambiente são modificados.

Os experimentos realizados não resolvem o problema geral da autonomia artificial. Ainda assim, demonstram que modificar a relação entre condição interna, percepção e curiosidade altera de maneira significativa o comportamento aprendido. Ao tornar a continuidade da percepção dependente das ações do agente, a precariedade permitiu que os recursos adquirissem relevância funcional sem uma recompensa direta para sua coleta. Esse resultado constitui um passo experimental em direção a agentes cujo comportamento seja orientado não apenas por objetivos externos, mas também pelas consequências de suas ações sobre suas próprias condições de interação com o mundo.